%% chapter3.tex — Methodology, Results, and Conclusions
%%
%% In a real thesis this would typically be split into separate files:
%%   \newchapter{chapter3}   % Methodology
%%   \newchapter{chapter4}   % Results and Discussion
%%   \newchapter{chapter5}   % Conclusions
%%
%% They are combined here to keep the template compact.

\chapter{Methodology, Results, and Conclusions}
\label{ch:methods}
\label{ch:conclusion}

%% -----------------------------------------------------------------------
%% 3.1 Methodology
%% -----------------------------------------------------------------------
\section{Methodology}
\label{sec:methodology}

Describe your experimental or analytical methods in enough detail that
a reader with the necessary background could reproduce your work.
Reference standard techniques with citations~\cite{kopka1999guide}
and reserve detailed derivations for appendices if they would interrupt
the narrative flow.

\subsection{Experimental Setup}
\label{subsec:setup}

Describe your equipment, software, datasets, and any parameters that
affect reproducibility.
If your setup involves multiple stages, a numbered list or a flow-diagram
figure helps the reader follow the sequence.

\subsection{Evaluation Metrics}
\label{subsec:metrics}

Define the metrics you use to evaluate your results.
Quantitative metrics should be accompanied by their mathematical
definitions, as in Equation~\eqref{eq:metric}:

\begin{equation}
  \label{eq:metric}
  \text{RMSE} = \sqrt{\frac{1}{N}\sum_{i=1}^{N}(y_i - \hat{y}_i)^2}.
\end{equation}

%% -----------------------------------------------------------------------
%% 3.2 Results
%% -----------------------------------------------------------------------
\section{Results and Discussion}
\label{sec:results}

Present your results objectively, then interpret them.
Connect each result back to the objectives stated in
Section~\ref{sec:objectives}.

When comparing to prior art (Table~\ref{tab:example} in
Chapter~\ref{ch:background}), be precise about the conditions under
which each method was evaluated.

\subsection{Quantitative Results}
\label{subsec:quantitative}

Summarise numerical results in tables, and use figures to reveal
trends or distributions that a table cannot capture clearly.

\subsection{Discussion}
\label{subsec:discussion}

Discuss the implications of your results.
What do they tell you about the research question posed in
Chapter~\ref{ch:introduction}?
Where do the results agree with prior work, and where do they differ?

%% -----------------------------------------------------------------------
%% 3.3 Conclusions
%% -----------------------------------------------------------------------
\section{Conclusions}
\label{sec:conclusions}

State the conclusions that can be drawn from your work.
Each conclusion should follow directly from a result or observation
presented in Section~\ref{sec:results}.

\begin{enumerate}
  \item The template compiles cleanly under TeX Live 2023 using
        \texttt{pdflatex} and \texttt{bibtex}.
  \item Per-chapter and single-list reference modes both function
        correctly with the provided \texttt{.bst} files.
  \item The preliminary pages (title, approval, abstract, TOC,
        List of Figures, List of Tables) are generated automatically.
\end{enumerate}

\section{Future Work}
\label{sec:future}

Identify open questions and directions for future investigation.
Future work might include:
\begin{itemize}
  \item Extension of the methodology to additional domains.
  \item Investigation of alternative approaches not explored here.
  \item Long-term evaluation studies.
\end{itemize}
